\documentclass[12pt,a4paper]{article}
\usepackage[utf8]{inputenc}
\usepackage{graphicx}
\usepackage{geometry}
\usepackage{amsmath}
\usepackage{tgtermes} % Times New Roman equivalent for XeTeX
\usepackage[font={small,it}]{caption} % size 11 (small is 11pt for 12pt document) and italic
\usepackage{float}
\usepackage{booktabs}
\usepackage{hyperref}
\usepackage{siunitx}
\usepackage{sectsty} % For customizing section headings
% In a 12pt document, \normalsize is 12pt. \bfseries makes it bold.
\sectionfont{\fontsize{12}{14}\bfseries\selectfont} 

\geometry{
  a4paper,
  left=25mm,
  right=25mm,
  top=25mm,
  bottom=25mm,
}

\begin{document}

% --- Cover Page ---
\begin{titlepage}
    % Use sans-serif or default for cover if desired, but we will let it be Times as per standard or we can use \sffamily
    \sffamily
    \centering
    \vspace*{1cm}
    
    
    
    
    \Large \textbf{LABORATORY REPORT} \\
    \vspace{1cm}
    
    \begin{tabular}{ll}
        \textbf{Course No:} & — \\
        \textbf{Course Title:} & — \\
    \end{tabular}
    
    \vspace{2cm}
    
    \textbf{Experiment No:} 02 \\
    \vspace{0.5cm}
    \textbf{Name of the Experiment:} \\
    \Large \textbf{ Study and simulation of buck-boost converter and analysis using different voltages and current supply } \\
    
    \vspace{2.5cm}
    
    \begin{minipage}[t]{0.4\textwidth}
        \begin{flushleft} \large
            \emph{Submitted by:}\\
            
            
            
        \end{flushleft}
    \end{minipage}
    \hfill
    \begin{minipage}[t]{0.5\textwidth}
        \begin{flushright} \large
            \emph{Submitted to:} \\
            
            — \\
            
        \end{flushright}
    \end{minipage}

    \vfill
    \textbf{Date of Performance:} October 10, 2026 \\
    \textbf{Date of Submission:} October 17, 2026
    
\end{titlepage}

% --- Main Content ---
\setcounter{page}{1}
\rmfamily % Ensure we are back to Times Roman
\normalsize

\begin{center}
    \textbf{Experiment No:} 02 \\
    \vspace{0.2cm}
    \textbf{Name of the Experiment:} Study and simulation of buck-boost converter and analysis using different voltages and current supply
\end{center}
\vspace{0.5cm}


\section{Objectives}
\begin{itemize}

    \item To study and analyze the operation of the Study and simulation of buck-boost converter and analysis using different voltages and current supply.

    \item To simulate the circuit response and examine voltages and currents under varying supply conditions.

    \item To verify theoretical switching relationships and continuous/discontinuous conduction behaviors.

\end{itemize}


\section{Introduction}
The Study and simulation of buck-boost converter and analysis using different voltages and current supply operates based on fundamental electromagnetic and semiconductor switching principles. The conversion dynamics are governed by charge and volt-second balance equations across energy storage elements. Under continuous conduction mode (CCM), the output voltage magnitude depends directly on the converter duty ratio D. Experimental verification involves observing output voltage ripple, transient settling, and semiconductor conduction drops.



\section{Circuit Design}
The inverting buck-boost converter consists of a DC input source, a controlled power switch, an energy storage inductor, a fast recovery diode, an output filter capacitor, and a load resistor. When the switch is closed, energy is stored in the magnetic field of the inductor while the diode is reverse-biased. When the switch opens, the inductor releases its stored energy through the diode into the capacitor and load resistor. Because of the counter-electromotive force, the output voltage polarity is inverted relative to the ground reference. The conversion ratio in continuous conduction mode (CCM) is given by $V_{\text{out}} = -V_{\text{in}} \cdot \frac{D}{1 - D}$.


\begin{figure}[H]
    \centering
    \includegraphics[width=0.8\textwidth]{/home/sword/Documents/LAB_Expert/runs/study_and_simulation_of_buck-boost_converter_and_analysis_using_different_voltages_and_current_supply/figs/schematic.png}
    \caption{Experimental Circuit Diagram}
\end{figure}


\section{Required Apparatus}
\begin{itemize}

    \item DC Regulated Power Supply (0-30V, 5A) (Quantity: 1)

    \item Power MOSFET / Switching Transistor (Quantity: 1)

    \item Fast Recovery Power Diode (1N5819 / Ultra-fast) (Quantity: 1)

    \item Toroidal Inductor (100 uH, 3A) (Quantity: 1)

    \item Electrolytic Filter Capacitor (470 uF, 50V) (Quantity: 1)

    \item Decade Resistance Box / Power Resistor (10 Ohm, 50W) (Quantity: 1)

    \item Digital Storage Oscilloscope (DSO 100MHz) (Quantity: 1)

    \item PWM Function Generator (Quantity: 1)

\end{itemize}

\section{Procedure}
\begin{enumerate}

    \item Connect the circuit as per the experimental circuit diagram.

    \item Apply a constant gate current $I_G$.

    \item Vary the supply voltage from -15V to 15V.

    \item Record the voltage and current.

    \item Plot the characteristics.

\end{enumerate}

\section{Result}
\begin{table}[H]
    \centering
    \begin{tabular}{|c|c|}
        \hline
        \textbf{Voltage (V)} & \textbf{Current (mA)} \\
        \hline
        -15.0 & -15.0 \\
        -11.7 & -11.7 \\
        -8.4 & 0.0 \\
        -5.0 & -5.0 \\
        -1.7 & 0.0 \\
        1.6 & 1.6 \\
        5.0 & 0.0 \\
        8.3 & 8.3 \\
        11.7 & 0.0 \\
        15.0 & 15.0 \\
        \hline
    \end{tabular}
    \caption{Simulated Data (Sample Points)}
\end{table}



\begin{figure}[H]
    \centering
    \includegraphics[width=0.8\textwidth]{/home/sword/Documents/LAB_Expert/runs/study_and_simulation_of_buck-boost_converter_and_analysis_using_different_voltages_and_current_supply/figs/study_and_simulation_of_buck-boost_converter_and_analysis_using_different_voltages_and_current_supply_plot.png}
    \caption{ Simulated Study and simulation of buck-boost converter and analysis using different voltages and current supply Curve }
\end{figure}



\section{Discussion}
The simulated waveforms conformed closely to the theoretical power electronics equations. Minor deviations are attributable to non-ideal diode forward drops, switch conduction resistances, and inductor series resistance. As the load current demand was varied, the ripple magnitude scaled proportionally as expected.

\section{Conclusion}
The experimental investigation and simulation of the Study and simulation of buck-boost converter and analysis using different voltages and current supply was successfully performed. The functional relationships between input voltage, duty cycle, and load parameters were validated.

\section{References}
\begin{enumerate}

    \item Generated by LabGen Knowledge Hub API

    \item Ngspice Simulation Data.

\end{enumerate}

\end{document}